% IF-10: explicit geometric coefficient realization, with every normalization retained.
\subsection{IF-10: a geometric realization over a CM elliptic curve}
There is a geometric realization of the finite coefficient comparison,
but its integral and character normalizations must be stated. It is not
an identification of an automorphic motive with the cohomology of one
abelian variety, and does not replace the operators of IF-08 by one $L$.

Let $K\subset\mathbb C$ be a number field containing the entries of the
finite matrices under consideration. Enlarge it to a number field
$K'$ containing $i$, $K$, and their complex conjugates, and write
$R=\mathcal O_{K'}$. Denominators of the matrices remain allowed and are
recorded by fractional $R$-lattices, as in (IF-08k). Use the elliptic curve
\[
 E:y^2=x^3-x\quad\text{over }\mathbb Q(i),\qquad
 [i](x,y)=(-x,iy).
\]
We work with algebraic correspondences modulo homological equivalence,
with coefficients in $K'$, and their Betti or $\ell$-adic realizations.
All correspondences constructed below are explicit combinations of
coordinate maps, the origin, $[i]$, and linear combinations of their
graphs. No Hodge-conjecture input is needed for this coefficient model.

\paragraph{Statement IF-10.}
For every finite chain of the exterior coefficient spaces and operators
in IF-08 whose normalized matrix entries lie in $K$, there is a chain
of homological motives generated by powers of $E$ and rank-one CM
character twists with the following properties.

\begin{enumerate}
\item Its realization has exactly the matrices of IF-08, including the
moment coefficients, determinant insertions, class sums, stabilizers,
local frames, invariant subspaces and terminal functionals. The
coefficient lattices, with their specified duals, transport isomorphically.
Before coefficient extension every Smith invariant of the same specified
map is identical. After extension, determinantal ideals are their extended
ideals; local valuations change by the ramification index, as specified
below. No equality of absolute finite indices across a field extension
is asserted.
\item Before the character twists, the polynomial coefficient space of
partition $\lambda$ is a submotive of a power of $E$ whose realization is
$T_\lambda\otimes M^{|\lambda|}$, where $M$ is the $i$-eigenmotive in
$h^1(E)$. A determinant twist $s$ changes the exponent to
$w=|\lambda|+rs$. The underlying integral polynomial lattice is a
saturated intersection in integral cohomology after extending coefficients
to $R$.
\item The Fock isometry uses $2^{-|\lambda|}$ times the ambient
Hermitian form on that polynomial eigenspace. The CM projector has
denominator $2$, and the dual eigenline introduces the factor $-2i$.
At primes above $2$ the transported abstract dual lattice must therefore
be distinguished from the native saturated conjugate eigenlattice.
No assertion of an integral cycle projector, or of an integral cup
polarization with the normalized Fock Gram matrix, is made.
\end{enumerate}

\begin{proof}
Choose an oriented integral symplectic basis $e,f$ of $H^1(E,\mathbb Z)$
such that
\[
 [i]^*e=-f,\qquad [i]^*f=e,\qquad
 \psi(e,f)=1.
\]
This is the basis $dx,dy$ on the square analytic torus. Thus
$u=e+if$ satisfies $[i]^*u=iu$ and
\[
 U=R u=\ker([i]^*-i)\subset H^1(E,R).
                                                               \tag{IF-10a}
\]
This line is saturated: $u,f$ form an $R$-basis, since their coefficient
matrix in $e,f$ has determinant one. Set
$\pi_0=[\{0\}\times E]$, $\pi_2=[E\times\{0\}]$, and
$\pi_1=\Delta_E-\pi_0-\pi_2$, with the convention that the displayed
correspondences act by pullback followed by pushforward to the second
factor. They project to cohomological degrees $0,2,1$, respectively.
The transpose graph of $[i]$ acts by $[i]^*$. Hence
\[
 p_+=\frac{1}{2}\bigl(\pi_1-i\,[i]^*\pi_1\bigr)
                                                               \tag{IF-10b}
\]
is an algebraic idempotent in the homological category and its image is
$M$. Its realization has image $K'u$ and its intersection with integral
cohomology is precisely $U$. The idempotence follows directly from
$([i]^*)^2=-1$ on $H^1$; homological equivalence makes this calculation
an equality of correspondences in the stated category.

For the usual Hermitian extension of the orthonormal form on $e,f$,
\[
 (u,u)_{\rm Herm}=2,\qquad
 \psi(u,\overline u)=-2i,\qquad
 (u,u)_{\rm bilinear}=0.                         \tag{IF-10c}
\]
These are three different pairings. Define the normalized Hermitian
form on $U$ by dividing its ambient Hermitian form by $2$; then $u$ has
norm one. Equivalently its value on $U$ is
$\psi(x,\overline y)/(-2i)$. This last expression is restricted to
$U$, where it is integral and positive under the chosen complex
embedding; it is not asserted to be an integral pairing on all of
$H^1(E,R)$.

For a column of height $h$, use the actual cohomological space
$H^h(E^r,K')$. For each $h$-element subset
$I=\{i_1<\cdots<i_h\}$ of the $r$ factors, project to $M$ in those
factors and to $h^0(E)$ in the others. Its integral generator is the
ordered cup product $u_{i_1}\cdots u_{i_h}$. These generators occupy
distinct K\"unneth summands and each is primitive, so their direct sum
is a saturated lattice canonically identified with
$\Lambda^h R^r\otimes U^{\otimes h}$. Applying this construction in
separate blocks for all column heights $h_c$ realizes
\[
 A_\lambda(R)=\bigotimes_c\Lambda^{h_c}R^r
 \quad\text{as}\quad
 A_\lambda(R)\otimes U^{\otimes d}
 \subset H^d(E^{rq},R),\qquad d=|\lambda|,
                                                               \tag{IF-10d}
\]
where $q$ is the number of columns. We use actual ordered cup products
and the resulting permutation signs. In particular this is the ordinary
exterior power of $H^1$ in cohomology, not an unqualified categorical
exterior power of the odd motive $h^1(E)$.

Every summand indexed by a tuple of subsets in (IF-10d) is isomorphic
to $M^{\otimes d}$. The isomorphism orders the selected factors and
discards the degree-zero factors. Its inverse inserts the origin and
the same degree-zero factors. These are induced by projections,
coordinate permutations, and the degree projectors. Compose the
projection onto one such summand, its ordered identification with
$M^{\otimes d}$, and the inverse identification for another summand.
The resulting algebraic correspondences $E_{JI}$ act by
\[
 E_{JI}(v_I)=v_J,\qquad E_{JI}(v_{I'})=0\quad(I'\ne I),
 \qquad E_{JI}E_{LK}=\delta_{I,L}E_{JK}.           \tag{IF-10e}
\]
The necessary signs are fixed by the chosen increasing order of the
degree-one factors. This provides all matrix units on the coefficient
space; it also provides rectangular matrix units between two ambient
powers with the same $d$.

Choose an integral basis for the primitive sublattice $T_\lambda$
proved in IF-08 and let $C$ be its matrix in the orthonormal basis of
$A_\lambda$. The rational matrix
\[
 P_T=C(C^tC)^{-1}C^t                              \tag{IF-10f}
\]
is the projector to its rational span. Inserting its entries into
(IF-10e) makes it an algebraic projector in the homological category.
Its integral intersection is $T_\lambda\otimes R$: primitivity over
$\mathbb Z$ gives a free quotient and is preserved by scalar extension.
All denominators in (IF-10f) are explicit divisors of
$\det(C^tC)$, before cancellation. Together with the factors $2$ from
(IF-10b), they record the denominator of this chosen correspondence.
The ambient tensor basis in (IF-10d) has squared norm $2^d$; after
multiplying the restricted Hermitian form by $2^{-d}$ its Gram matrix
on $T_\lambda$ is the one in IF-04. Thus the normalized Fock isometry
is exact, including its pairing normalization.

Every coefficient matrix $A_\beta$ of IF-08, and hence every
$B_J=\sum_\beta J(\beta)A_\beta$ with $J(\beta)\in K'$, is now an
algebraic correspondence by (IF-10e), restricted on each side by
(IF-10f). Its matrix on the chosen lattice is unchanged. This argument
does not assert that the coefficients of the representing cycle are
integral; the matrix on the specified lattice is integral exactly
when the corresponding matrix of IF-08 is integral.

To handle determinant characters and arrows of different total
degrees, note that the determinant line of $U^{\oplus r}$ is
$M^{\otimes r}$ with generator $u^{\otimes r}$, using ordered cup
products. Negative tensor powers use the abstract dual of the chosen integral
line from the outset, including in the determinant twist. The polarization
identifies $M^\vee$ with the conjugate eigenmotive $\overline M(1)$;
the scalar in (IF-10c) specifies its integral normalization. These
rank-one motives are invertible in the homological category: evaluation
and its transpose, divided by their nonzero scalar, are inverse on the
one-dimensional realization, hence inverse there. Consequently the
coefficient motive for $(\lambda,s)$ is
\[
 \mathcal V_{\lambda,s}=T_{\lambda,K'}\otimes M^w,
 \qquad w=d+rs.
\]
Normalize it by the specified inverse CM character:
\[
 \mathcal V_{\lambda,s}\otimes M^{-w}
       \simeq T_{\lambda,K'}\otimes\mathbf1.       \tag{IF-10g}
\]
Use the abstract dual of the chosen integral rank-one lattice for this
evaluation, so that its generator evaluates to one. Formula (IF-10c)
says exactly what this means inside a conjugate geometric realization:
$U^\vee$ is represented by
$(-2i)^{-1}\overline U(1)$. In particular it is not the native
saturated lattice $\overline U(1)$ at primes above $2$. If one instead
uses native integral eigenlines in a pair of opposite powers
$M^w,M^{-w}$, evaluation has image the ideal generated by
$(2i)^{|w|}$, up to a unit and the ordering sign. This factor is retained,
or is removed by precisely the displayed abstract-dual normalization;
there is no unnoticed saturation at $2$. This $|w|$ formula concerns
the already collapsed pair $M^w,M^{-w}$. If $s<0$ and one begins
instead with the raw native tensor of $d$ copies of $U$ and
$r|s|$ copies of $\overline U(1)$, converting the latter to abstract
duals contributes $(-2i)^{-r|s|}$. Collapsing the raw native tensor by
pairings first contributes $(-2i)^{\min(d,r|s|)}$, up to the ordering
unit. These factors must be retained in any comparison with that
different native presentation.

After (IF-10g), every rectangular matrix, including the determinant
insertion and maps between different $w$, is an algebraic correspondence
between sums of the unit motive. This realizes exactly the original
coefficient matrices because the evaluation generators were fixed to
one. Finite direct sums, invariant subspaces and level subspaces are
handled by matrix projectors over $K'$ and their specified integral
intersections. An arbitrary projective global lattice is retained as
that lattice; a free global basis is not assumed. Thus the block sums
and local-frame conjugations of (IF-08k), and the terminal dual
functional, are unchanged.

Finally a finite torsion-free extension of rings of integers is
faithfully flat. The minors defining every determinantal ideal are
the same minors after extension to $R$; contracting the extended ideal
recovers the original one. Over a local extension of discrete valuation
rings, a valuation is multiplied by the ramification index, and this
known factor recovers the original length. This proves the index
assertion at every prime, with the abstract-dual and cycle-denominator
conventions just specified.
\end{proof}

\paragraph{The actual algebraicity input.}
The coefficient-field hypothesis is already part of a supplied
algebraic theta-chain datum, and is not proved by calling an arbitrary
complex derivative algebraic. The exact interfaces are
\texttt{ThDistC.lean}, \texttt{Amb.JetsRat} (lines 240--242), and
\texttt{KernelLocalFactors/ThDistC.lean}, \texttt{jetK\_spec}
(lines 615--621), which takes that rationality hypothesis and identifies
the descended moment with the period-normalized derivative.
For the actual tied kernel, \texttt{JetFieldLeKR.lean},
\texttt{KerEntTied.KerJets} (lines 113--116), includes the product
\texttt{rhoDet*rhoDet'} in the algebraic coefficient. One therefore
uses the combined moment $\widetilde J=\rho\rho'J$ when only that
combination is certified algebraic; algebraicity of its separate
factors is unnecessary. Alternatively,
\texttt{DataForms.lean}, \texttt{KRat,KL,KL\_spec} (lines 252--259),
gives the algebraic global map together with its exact complex
realization. These declarations are hypotheses and conditional descent
theorems; they do not by themselves construct a new closed global
instance. IF-10 applies to every instance supplied with this data. In
particular the fixed \texttt{ThE} test of IF-09 is unconditional over
$\mathbb Q$ at the level of its finite moments.

\paragraph{Actual weights and the limit of Tate normalization.}
For one embedding factor $\tau$ of the three actual column coefficients
in IF-09, direct summation gives
\[
 (|\alpha_0|,|\alpha_1|,|\alpha_2|)=(113,841,889),
 \qquad (w_0,w_1,w_2)=(137,-1463,-3351).           \tag{IF-10h}
\]
At the modified real place of \texttt{numBig}, all three $w_j$ decrease
by $3$; their successive differences remain $-1600,-1888$.
The actual coefficient module is a tensor product over $\tau$; apply
(IF-10g) to each factor and then tensor. If the same $M$ is used at
every place, its total character exponent is $\sum_\tau w_{j,\tau}$,
not a single entry of (IF-10h). The CM twists in (IF-10g) are therefore
substantial and are not pure
Tate twists. Indeed, under the selected embedding $M^w$ has Hodge type
$(w,0)$. A nonzero algebraic correspondence between
$M^w(a)$ and $M^{w'}(b)$ would preserve Hodge type, requiring
\[
 (w-a,-a)=(w'-b,-b).
\]
It follows that $a=b$ and $w=w'$. Hence in this natural single-eigenline
model the nonzero arrows of the actual coefficient chain cannot become
algebraic by Tate twists alone. The explicit CM-character normalization
is necessary for this construction. This is an obstruction to this
specified model, not a claim about every conceivable realization.

On an $\ell$-adic realization, if the eigenline has character $\eta$,
the untwisted coefficient motive has character $\eta^w$ on its rank-one
factor (with the sum of exponents when tensoring over $\tau$). Formula
(IF-10g) cancels that character. Thus this construction
does not preserve an unspecified automorphic $L$-function or supply its
motivic realization. It does give algebraic correspondences realizing
the fully specified finite coefficient matrices and their exact local
indices. In the language of cote 16, scans 014--018, it constructs
operators in a geometric correspondence ring after the stated
normalizations; IF-09 prevents replacing all such two-step kernels by
positive powers of one fixed Lefschetz element. The $\lambda$-operations
occur in their actual cohomological exterior powers, while $K_0$ alone
forgets the retained lattice indices and the prime-$2$ dual-lattice
factor.

The equation for $E$ and these coordinate correspondences spread over
$\mathbb Z[i,1/2]$. Statements about the transported lattice at rational
prime $2$ concern integral Betti cohomology or the specified $2$-adic
representation lattice. They do not assert good reduction of this
elliptic curve at $2$, or an integral Chow projector there.
